% !PRINTED 254
\weberpage{254}

\begin{center}
\S.\,66.\\
\gesperrt{Einfachheit der Gruppe $E$.}
\end{center}

Die erste Frage bei einer eingehenderen Untersuchung der Gruppe $E$ ist die nach einem etwa vorhandenen Normaltheiler. Es l\"asst sich beweisen, dass ein solcher Normaltheiler, von zwei ganz einfachen Ausnahmen abgesehen, nicht vorhanden ist.

Nehmen wir an, es sei $G$ ein Normaltheiler von $E$, der wenigstens eine von der Identit\"at verschiedene Substitution
\weq{1}{A = \begin{pmatrix} \alpha, & \beta \\ \gamma, & \delta \end{pmatrix}}
enth\"alt. Nach dem Begriffe des Normaltheilers ist dann, wenn $T$ ein beliebiges Element in $E$ ist,
\weq{2}{TAT^{-1}}
auch in $G$ enthalten, und es ist zu zeigen, dass man auf diese Weise und durch Zusammensetzung solcher Substitutionen alle Elemente von $E$ ableiten kann, woraus dann folgt, dass $G$ mit $E$ identisch sein muss. Es gen\"ugt aber dazu, das Vorkommen der erzeugenden Substitutionen $A_0$, $B_h$ [\S.\,65, (8)] in $G$ nachzuweisen.

Wir gehen dazu schrittweise vor.

1) In $G$ kommt eine Substitution von der Form
\weq{3}{\begin{pmatrix} 0, & \beta \\ \gamma, & \delta \end{pmatrix}}
vor. Um dies zu zeigen, bilden wir nach (2):
\weq{4}{\begin{pmatrix} \xi, & 1 \\ -1, & 0 \end{pmatrix}\begin{pmatrix} \alpha, & \beta \\ \gamma, & \delta \end{pmatrix}\begin{pmatrix} 0, & -1 \\ 1, & \xi \end{pmatrix} = \begin{pmatrix} \beta\xi+\delta, & \beta\xi^2+(\delta-\alpha)\xi-\gamma \\ -\beta, & -\beta\xi+\alpha \end{pmatrix}.}
Wenn nun $\beta$ von Null verschieden ist, so kann man $\xi$ aus $\beta\xi+\delta=0$ bestimmen und hat das Ziel erreicht. Ist aber $\beta=0$, so ist nicht gleichzeitig $\gamma=0$, $\delta-\alpha=0$, weil sonst $A$ die identische Substitution w\"are; man kann daher $\xi$ so annehmen, dass $(\delta-\alpha)\xi-\gamma$ von Null verschieden ist. Dann hat man durch (4) eine Substitution $A$ in $G$ gebildet, in der $\beta$ nicht Null ist, und auf die man die Formel (4) nun nochmals anwenden kann, um $\alpha$ zu Null zu machen.

2) In $G$ kommt eine Substitution von der Form
\[
\begin{pmatrix} 0, & 1 \\ -1, & \xi \end{pmatrix}
\]
%CONT>>
